\documentclass[10pt,a4paper]{amsart}
\usepackage[margin=19mm]{geometry}
\usepackage{amsmath,amssymb,mathtools}
\usepackage[scr=rsfs,cal=euler]{mathalfa}
\usepackage{xfrac}
\usepackage[unicode,hidelinks,bookmarks=false]{hyperref}
\newcommand{\ie}{i.e.\ }
\newcommand{\cf}{cf.\ }
\newcommand{\resp}{respectively\ }
\newcommand{\Subsection}{\subsection}
\providecommand{\nobreakdash}{\nobreak\hbox{-}}
\setlength{\emergencystretch}{2em}
\multlinegap0pt
\usepackage{amssymb}
\usepackage{tikz}
\usepackage{mathtools}
\usepackage[shortlabels]{enumitem}
\newtheorem{lemma}{Lemma}
\title{$K$-Theoretic Obstructions to Linearizing QCA Representations}
\author{Mattie Ji, Bowen Yang}
\date{Source: arXiv:2606.19657v1 (CC BY 4.0)}
\begin{document}
\maketitle

\noindent\emph{Source and licence.} $K$-Theoretic Obstructions to Linearizing QCA Representations. Version arXiv:2606.19657v1 (CC BY 4.0), distributed by arXiv under the Creative Commons Attribution 4.0 International licence, http://creativecommons.org/licenses/by/4.0/, as recorded in the arXiv licence metadata for this exact version and shown on the abstract page https://arxiv.org/abs/2606.19657v1. Full licence notice: \texttt{licenses/arxiv-2606.19657-CC-BY-4.0.txt}.\par\smallskip

\noindent\textit{Editorial scope.} Counted unit: the article's lemma lem:homological-killing (source0.tex lines 1181--1183). The article's complete original proof of it is delivered with it (source0.tex lines 1290--1331, one \texttt{proof} environment of 42 lines, which the article places away from the statement). No mathematics is written, repaired or re-derived: the statement and the proof are the article's own lines, in the article's own notation, and no retained line is substituted or re-flowed.  No further source-local context is needed: every cross-reference of every retained line resolves inside the two delivered spans.  Nothing was omitted from any retained span, so the package carries no omission record.  No retained line contains a citation, so the wrapper carries no bibliography and no bibliography entry.  No retained line contains a figure environment, an image-inclusion command or drawing code, so no image is delivered and the package declares no asset dependency.  Editorial formatting only, in addition: an AMS \texttt{amsart} preamble that restates the article's own declarations and macros used by the retained lines, namely the environments \texttt{lemma} on separate counters, together with the standard packages those lines need.  The retained environments print in this document as Lemma 1 in order of appearance, whereas the article prints the same items with the numbering of its own full document.  The complete native source of the article as distributed in the arXiv e-print for this exact version is bundled unchanged as \texttt{sources/203160/source0.tex}.
\begin{lemma} \label{lem:homological-killing}
The space $X_{n+1}$ is homologically $(n+1)$-connected.
\end{lemma}

\begin{proof}[Proof of Lemma~\ref{lem:homological-killing}]
Consider the Serre spectral sequence of the fibration
\begin{equation}
K(A_n,n)\longrightarrow X_{n+1}\longrightarrow X_n .
\end{equation} with
\begin{equation}
E^2_{p,q}\cong H_p(X_n;H_q(K(A_n,n);\mathbb Z)).
\end{equation}

The fiber $K(A_n,n)$ is $(n-1)$-connected. In low degrees, its homology is 
\begin{equation}
H_q(K(A_n,n);\mathbb Z)\cong
\begin{cases}
\mathbb Z, & q=0,\\
0, & 0<q<n,\\
A_n, & q=n,\\
0, & q=n+1.
\end{cases}
\end{equation} On the other hand, $H_p(X_n;\mathbb Z)=0$ for $0<p\leq n$.

It follows that, in total degree at most $n+1$, the only possibly nonzero terms are $E^2_{0,0}\cong \mathbb Z$, $E^2_{n+1,0}\cong H_{n+1}(X_n;\mathbb Z)$, and $E^2_{0,n}=H_0(X_n, H_n(K(A_n,n);\mathbb Z))$.

The differential has bidegree
\begin{equation}
d^r:E^r_{p,q}\longrightarrow E^r_{p-r,q+r-1}.
\end{equation}
The term $E^r_{n+1,0}$ receives no differentials, since an incoming differential would have source $E^r_{n+1+r,1-r}$, which vanishes for $r\geq 2$. 

For $2\leq r\leq n$, the outgoing differential
\begin{equation}
d^r:E^r_{n+1,0}\longrightarrow E^r_{n+1-r,r-1}
\end{equation}
has target in a row $0<r-1<n$, and this row vanishes because $H_{r-1}(K(A_n,n);\mathbb Z)$ vanishes. Therefore the first possible nonzero differential out of $E_{n+1,0}$ is $d^{n+1}.$
Under the identifications above, this is a homomorphism 
\begin{equation}
d^{n+1}:H_{n+1}(X_n;\mathbb Z)\to H_0(X_n, H_n(K(A_n,n);\mathbb Z))\cong A_n.
\end{equation}

We now identify this differential. By naturality of the Serre spectral sequence the transgression $d^{n+1}$ is evaluation against the classifying class $\iota_n\in H^{n+1}(X_n;A_n)$. Therefore the differential is an isomorphism.

Consequently, both terms $E^{n+1}_{n+1,0}$ and $E^{n+1}_{0,n}$ are killed by this differential. It follows that $\widetilde H_i(X_{n+1};\mathbb Z)=0$ for $i\leq n+1$. Therefore $X_{n+1}$ is homologically $(n+1)$-connected.
\end{proof}

\end{document}
